\section{Powder Synthesis and Precursor Chemistry}
\label{sec:powder}

Powder synthesis sets the chemical length scale that the later firing step must eliminate. Solid-state routes are robust and scalable, but their reaction path depends on oxide surface area, mixing order and lithium precursor decomposition. Sol--gel, combustion, spray-flame and mechanochemical routes shorten diffusion distances or increase homogeneity, yet they introduce organics, residual carbon, agglomeration or a different lithium-loss profile \cite{puente2021garnet,ali2021spray,alizadeh2023synthesis,duvel2012mechanosynthesis,kalita2023recent}.

\subsection{Route families}

Conventional oxide or carbonate mixing remains the reference route because precursor identities and mass balances are easy to audit. A comparative Al--LLZO study found that solid-state and sol--gel routes can yield different thermal, structural and electrochemical responses even at nominally similar composition \cite{kosir2022comparative}. Reaction-pathway-informed synthesis further shows that a short calcination can be effective when precursor availability and intermediate phases are controlled, rather than when dwell time is simply reduced \cite{garman2025reaction}. These findings favor a route description that records milling media, order of addition, solvent or binder, calcination atmosphere, intermediate phase and the reason for selecting the dwell time.

Spray and combustion approaches can generate fine or chemically mixed particles at lower apparent synthesis temperatures \cite{ali2021spray,alizadeh2023synthesis,djenadic2014nebulized}. Mechanochemical processing similarly changes defect density and contact area before crystallization \cite{duvel2012mechanosynthesis}. The benefit is not universal: high surface area accelerates both reaction and environmental degradation, and agglomerates can create macropores that survive pressing. Particle-size-dependent degradation in Ta-containing garnet powders makes air exposure time and storage atmosphere essential metadata \cite{hoinkis2023particle}.

The route family is broader than the three commonly compared labels. Low-temperature synthesis reviews, spray-flame work, high-entropy garnet studies and AI-guided precursor screening each target shorter diffusion distances or improved compositional uniformity, but they also change agglomeration, defect populations and lithium-loss kinetics \cite{kalita2023recent,ali2021spray,feng2023discovery,wang2026ai,kim2025machine,li2026green,li2025phase}. Processing guidelines therefore recommend reporting powder morphology, precursor hydration, atmosphere and the thermal event that removes organics or carbon \cite{balaish2025emerging,puente2021garnet}.

Low-temperature Bi substitution and carbonate crystallization kinetics show why powder handling belongs in the synthesis record: defect chemistry and CO$_2$/LiOH exposure can alter the surface before the garnet reaction is complete \cite{schwanz2019ionic,su2026kinetics,huo2019li}.

\subsection{Lithium precursors and mass balance}

Li$_2$CO$_3$ is convenient but can leave carbonate contamination or volatilize during high-temperature treatment; LiOH changes decomposition and water release. Recent precursor comparisons therefore treat lithium source as a process variable rather than a procurement detail \cite{bae2026effect,shi2025li2co3,touidjine2024the}. A reproducible batch record should state the assay-corrected lithium content, intentional excess, hydration state, powder-bed cover and the measured mass before and after calcination. When these quantities are unavailable, the phrase “lithium compensated” should be read as nominal rather than measured.

Direct precursor comparisons support this distinction. Li$_2$CO$_3$ and LiOH produce different decomposition paths in LLZNO and LLZTO, and Ta-containing batches can retain a composition that differs from the weighed recipe after firing \cite{bae2026effect,li2025synthesis,shi2025li2co3,liu2014excess,narayanan2015effect}. Older Al-, W- and Ge-containing synthesis studies likewise show that calcination, milling and the choice of complexing route set the intermediate phases before densification \cite{jin2011al,li2015w,rosenkiewitz2014nitrogen,chen2014effect,kokal2011sol}.

\subsection{From powder to green body}

Pressing pressure, binder removal and particle-size distribution control the green-body pore network that later governs shrinkage. A small powder that reacts early may densify differently from a coarse powder with the same phase composition. The practical comparison is thus a matched green density and thickness, followed by a thermal profile that records shrinkage or dilatometry. Low-temperature synthesis studies and scalable tape-casting reviews emphasize that powder morphology and slurry rheology must be optimized together with the final membrane thickness \cite{kalita2023recent,touidjine2026tapex,gao2019preparation}.

For transfer to thin or composite forms, the relevant precedent includes submicron Nb-doped powders, rapid Li compensation in films and tape-cast garnet sheets \cite{yan2020submicron,xiang2023rapid,jonson2018tape,teng2013recent}. These studies motivate matching green density and thickness before comparing the later thermal schedule, rather than ranking powders by BET area alone.

\begin{table}[t]
\centering
\caption{Powder-route comparison with the variables that determine transferability.}
\label{tab:powder}
\begin{tabular}{P{0.20\linewidth}P{0.24\linewidth}P{0.25\linewidth}P{0.20\linewidth}}
\toprule
Route family & Main opportunity & Dominant uncontrolled risk & Minimum report \\
\midrule
Oxide/carbonate solid state & Simple scale-up and explicit mass balance & Incomplete mixing and long diffusion path & Milling sequence, particle size, calcination and Li excess \cite{heywood2023tailoring,kosir2022comparative} \\
Sol--gel or combustion & Chemical homogeneity and fine particles & Organics, residual carbon and agglomeration & Precursor concentration, burnout and powder morphology \cite{alizadeh2023synthesis,djenadic2014nebulized} \\
Spray or mechanochemical & Short diffusion distance and defect engineering & Rapid surface reaction, broad size distribution & Residence or milling history, atmosphere and storage \cite{ali2021spray,duvel2012mechanosynthesis,hoinkis2023particle} \\
\bottomrule
\end{tabular}
\end{table}
